%%%%%%%%%%%%%%%%%%%%%%%%%%%%%%%%%%%%%%%%%%%%%%%%%%%%%%%%%%%%%%%%%%%%
%% thesis.tex
%% Master Thesis Master Document
%% University of Europe for Applied Sciences (UE), Potsdam
%% Author: Watcharapon Anaketanaset (Matriculation No. 94658335)
%% Supervisor: Prof. Dr. Iftikhar Ahmed
%%%%%%%%%%%%%%%%%%%%%%%%%%%%%%%%%%%%%%%%%%%%%%%%%%%%%%%%%%%%%%%%%%%%
\documentclass{uethesis}

\usepackage{booktabs}
\usepackage{amsmath}
\usepackage{amssymb}
\usepackage{array}

%%%%%%%%%%%%%%%%%%%%%%%%%%%%%%%%%%%%%%%%%%%%%%%%%%%%%%%%%%%%%%%%%%%%
%% Set Thesis Information (Approved Proposal & UE Potsdam Directives)
%%%%%%%%%%%%%%%%%%%%%%%%%%%%%%%%%%%%%%%%%%%%%%%%%%%%%%%%%%%%%%%%%%%%
\renewcommand{\thesistitle}{Does National Specialisation Help Forecast Research-Topic Growth? Evidence from Germany, China, and Japan Using OpenAlex}
\renewcommand{\thesisauthor}{Watcharapon Anaketanaset}
\renewcommand{\matriculation}{94658335}
\renewcommand{\firstsupervisor}{Prof. Dr. Iftikhar Ahmed}
\renewcommand{\secondsupervisor}{Dr. Shan Faiz}
\renewcommand{\submissiondate}{February 12, 2027}
\renewcommand{\department}{Department of Tech and Software Engineering}
\renewcommand{\degree}{Master of Science in Data Science (120 ECTS)}

%%%%%%%%%%%%%%%%%%%%%%%%%%%%%%%%%%%%%%%%%%%%%%%%%%%%%%%%%%%%%%%%%%%%
%% Define Abstract Content
%%%%%%%%%%%%%%%%%%%%%%%%%%%%%%%%%%%%%%%%%%%%%%%%%%%%%%%%%%%%%%%%%%%%
\renewcommand{\theabstract}{
Anticipating which scientific and technological topics will experience rapid expansion is of critical importance for researchers, research institutions, and national science policy makers. While prior research forecasting models have predominantly relied on past aggregate global publication history, national research systems frequently serve as concentrated incubators for specialized technologies before global diffusion occurs. This thesis examines whether incorporating national publication volumes and national scientific specialisation---measured via Fractional All-Science Revealed Comparative Advantage ($RCA_{\text{AllSci}}^{\text{Frac}}$)---improves the forecast accuracy of future 3-year global research-topic growth across 1,025 technology topics indexed in OpenAlex from 2010 to 2024.

Focusing on three major scientific powers with contrasting institutional architectures---Germany (DE), China (CN), and Japan (JP)---the study constructs a leakage-free, balanced longitudinal panel encompassing 15,375 topic-year observations restricted to selected scholarly work types (`article`, `conference-paper`, `review`). Publications are attributed using server-side equal fractional counting ($1/k$), eliminating multi-country co-authorship inflation while preserving census coverage. A nested hierarchy of regularized linear models ($M0$ to $M3$) is evaluated across two strictly non-overlapping temporal test horizons (Test Origins 2018 and 2021). Both continuous forecast accuracy (Mean Absolute Error) and early identification of emerging high-growth topics outside the historical top growth quintile (Precision@50) are assessed with bootstrap confidence intervals. The findings provide rigorous empirical evidence regarding the predictive utility of national specialization in technology trend forecasting.

\vspace{0.5cm}
\noindent \textbf{Keywords:} Research Trend Forecasting, Scientometrics, OpenAlex, National Specialisation, Revealed Comparative Advantage, Fractional Counting, Emerging Technologies
}

%%%%%%%%%%%%%%%%%%%%%%%%%%%%%%%%%%%%%%%%%%%%%%%%%%%%%%%%%%%%%%%%%%%%
%% Begin Document
%%%%%%%%%%%%%%%%%%%%%%%%%%%%%%%%%%%%%%%%%%%%%%%%%%%%%%%%%%%%%%%%%%%%
\begin{document}

% Front matter: Title page, Author Declaration, Abstract, Table of Contents, Lists of Tables & Figures
\beginthesis

%%%%%%%%%%%%%%%%%%%%%%%%%%%%%%%%%%%%%%%%%%%%%%%%%%%%%%%%%%%%%%%%%%%%
%% Main Chapters
%%%%%%%%%%%%%%%%%%%%%%%%%%%%%%%%%%%%%%%%%%%%%%%%%%%%%%%%%%%%%%%%%%%%
\chapter{Introduction} \label{chap:intro}

\section{Background}
Scientific research and technological development progress at an uneven pace: while certain scientific topics experience exponential expansion and redefine technological paradigms, others grow modestly, stagnate, or diminish over time. The capacity to anticipate which research domains will undergo accelerated growth holds profound implications for researchers, corporate strategists, academic institutions, and national science policy makers \citep{rotolo2015}. Governments worldwide allocate substantial portions of public funding to identify and foster ``critical,'' ``emerging,'' or ``strategic'' technologies. However, strategic allocation requires robust, timely, and objective forecasting mechanisms that can navigate the inherent uncertainty of nascent scientific domains.

The concept of an ``emerging technology'' encompasses multiple qualitative and quantitative dimensions, including radical novelty, relatively fast growth, coherence, prominent socio-economic impact, and non-negligible uncertainty \citep{rotolo2015}. Among these dimensions, publication growth constitutes a central empirical trace through which scientific emergence manifests and can be systematically monitored. Prior computational bibliometric studies have leveraged publication counts, citation trajectories, and machine learning models to forecast topic expansion \citep{liang2021}. More recently, scholarship in science dynamics and economic complexity has demonstrated that national scientific specialisation reflects institutional capabilities and predictive path dependencies within national research spaces \citep{galanmena2026}. 

This thesis investigates a distinct and unexplored question: \emph{Can national scientific specialisation help forecast how fast a technology-related research topic will grow globally, and can it assist in identifying topics that have not yet achieved high growth at the forecast date?}

To investigate this question empirically, this study examines three major national research systems with contrasting institutional architectures, geopolitical priorities, and developmental trajectories:
\begin{enumerate}
    \item \textbf{Germany (DE):} A mature, highly internationalised, and decentralised European science system characterized by long-standing strengths in precision engineering, materials science, and chemical engineering.
    \item \textbf{China (CN):} An rapidly expanding, state-directed scientific ecosystem driven by strategic technological modernization, whose national research and development (R\&D) expenditure increased substantially to approximately 2.7\% of gross domestic product (GDP) in 2024 \citep{nbs2025}.
    \item \textbf{Japan (JP):} A technologically advanced East Asian research system with deep historical capabilities in hardware, energy, and robotics, operating under demographic transition and structural R\&D restructuring.
\end{enumerate}
Together, these three countries represent a major portion of global technological output and provide a rich comparative testbed for evaluating whether national specialization patterns encode predictive signals regarding global topic dynamics.

\section{Problem Statement}
Existing approaches to research trend forecasting predominantly rely on global-level indicators, such as past aggregate publication volume, citation velocity, and textual term burstiness \citep{liang2021}. While global historical momentum provides a robust baseline, it suffers from a fundamental limitation: by definition, global aggregate counts are lag indicators that reflect consensus that has already occurred across the scientific community. By the time a topic achieves high global volume, it is no longer nascent.

Conversely, individual national research systems frequently act as early incubators or concentrated testbeds for specific technological trajectories before widespread global diffusion takes place. For example, national technology roadmaps, concentrated industrial consortia, or state-funded target programs in nations such as China, Germany, or Japan may mobilize substantial research activity into specialized niches years prior to broader international adoption. 

However, whether national scientific activity genuinely provides advance forecasting power beyond naive momentum and global history remains an open empirical question. Specifically:
\begin{itemize}
    \item Does observing publication volume from individual leading countries add predictive accuracy over global publication history alone?
    \item Does relative national specialization---measured via Revealed Comparative Advantage (RCA)---provide informative signal beyond mere absolute national publication counts?
    \item Can national specialization assist in the early identification of emerging topics---specifically identifying topics that are currently outside the top-growth tier but will subsequently experience high growth?
\end{itemize}
Addressing these questions requires addressing significant methodological challenges: eliminating lookahead bias, enforcing leakage-free temporal training and testing boundaries, properly accounting for multi-country co-authorship via fractional attribution, and establishing transparent, parameter-free momentum baselines.

\section{Research Objectives}
The primary objective of this thesis is to empirically determine whether incorporating national publication volumes and fractional Revealed Comparative Advantage (RCA) from Germany, China, and Japan improves the forecasting of future three-year research-topic growth rate across 1,025 technology topics indexed in OpenAlex \citep{priem2022}.

The specific objectives of the study are:
\begin{enumerate}
    \item To construct a leakage-free, balanced longitudinal panel spanning 15 annual intervals (2010--2024) across 1,025 physical science and technology topics, incorporating fractional publication counts and all-science national denominators.
    \item To analyze the longitudinal evolution of technology topics and national scientific specialization across Germany, China, and Japan over the 2010--2024 period.
    \item To build and evaluate a nested hierarchy of regularized linear forecasting models ($M0$ to $M3$) across two strictly non-overlapping temporal test horizons (Test Origin 2018: 2019--2021; Test Origin 2021: 2022--2024).
    \item To determine whether national specialization adds incremental predictive power beyond absolute national publication counts and global publication history.
    \item To assess the utility of national specialization in identifying emerging, high-growth topics early, specifically focusing on topics not previously in the top growth quintile.
\end{enumerate}

\section{Research Questions}
To achieve these objectives, the thesis addresses three formal Research Questions (RQs) approved in the research proposal:
\begin{itemize}
    \item[\textbf{RQ1:}] \emph{How have technology-related research topics and national specialisation developed over time in Germany, China, and Japan?}
    \item[\textbf{RQ2:}] \emph{Do country publication counts and specialisation improve forecasts of topic growth beyond global publication history? Does specialisation add anything beyond counts?}
    \item[\textbf{RQ3:}] \emph{Do country counts and specialisation help identify topics that become high-growth but were not high-growth at the forecast date?}
\end{itemize}

\section{Contributions}
This thesis provides several distinct contributions to the fields of scientometrics, science and technology policy, and quantitative research forecasting:
\begin{itemize}
    \item \textbf{Methodological Rigor and Leakage Elimination:} Unlike studies that employ retrospective cross-validation across randomly partitioned topics, this research enforces strict origin-specific eligibility criteria ($P_{[t-2, t]} \ge 50$ works) and temporal train/test partitioning, completely eliminating lookahead bias and outcome leakage.
    \item \textbf{Equal Fractional Country Attribution ($1/k$):} In response to supervisory guidance and bibliometric best practice, the study implements server-side equal fractional counting for international collaborative works, eliminating the inflationary double-counting bias inherent in whole counting while preserving 100\% group census coverage.
    \item \textbf{Aligned Global Normalization:} The target variable and the naive momentum baseline are mathematically aligned by explicitly subtracting global all-science growth under the identical selected scholarly work types scope (`article|conference-paper|review`), isolating true technological topic growth from the underlying expansion of the bibliographic database.
    \item \textbf{High-Growth Early Identification Benchmark (RQ3):} The study operationalizes a rigorous decision-theoretic evaluation protocol for emerging topic identification by conditioning specifically on topics outside the top growth quintile at the forecast origin, testing model precision against a formal 20\% chance level.
    \item \textbf{Open Science and Reproducibility:} The complete master dataset, synchronized denominator tables, and deterministic extraction and audit pipelines are permanently archived with cryptographic checksums (SHA-256) in a public repository.
\end{itemize}

\section{Thesis Organization}
The remainder of this thesis is structured as follows:
\begin{itemize}
    \item \textbf{Chapter~\ref{chap:background} (Background):} Introduces the theoretical foundations of emerging technologies, scientometric indicators, economic complexity, and the mathematical definition of Revealed Comparative Advantage (RCA).
    \item \textbf{Chapter~\ref{chap:lit} (Literature Review):} Critically reviews prior literature on research trend forecasting, bibliometric predictors, national scientific diversification, and identifies the core research gap addressed by this work.
    \item \textbf{Chapter~\ref{chap:method} (Methodology):} Details the OpenAlex data extraction, selected scholarly work types filtering, fractional attribution mechanism, dataset synchronization, feature engineering hierarchy ($M0$--$M3$), temporal evaluation splits, and bootstrap inference protocols.
    \item \textbf{Chapter~\ref{chap:res} (Results and Discussions):} Presents empirical findings for RQ1 (historical specialization trajectories), RQ2 (forecast accuracy across model tiers), and RQ3 (early identification of emerging high-growth topics), accompanied by in-depth discussion and sensitivity analyses.
    \item \textbf{Chapter~\ref{chap:conc} (Conclusions):} Summarizes the primary findings, discusses practical implications for science policy and R\&D investment, explicitly acknowledges methodological limitations, and outlines fruitful avenues for future research.
\end{itemize}

\chapter{Background} \label{chap:background}

\section{Emerging Research Topics}
% [Insert text on defining emerging technologies, e.g., novelty, growth, impact, uncertainty (Rotolo et al., 2015)]

\section{Bibliometrics}
% [Insert text on bibliometric data, scholarly publication indicators, and measurement concepts]

\section{National Specialisation}
% [Insert text on national research systems, capabilities, and country-level diversification]

\section{Revealed Comparative Advantage (RCA)}
% [Insert mathematical definition, Balassa formula, and application to scientometrics]

\section{Research Forecasting Concepts}
% [Insert overview of forecasting principles, momentum, and prediction horizons]

\chapter{Literature Review} \label{chap:lit}

\section{Previous Studies in Research Topic Forecasting}
Forecasting emerging scientific and technological topics has attracted significant interest across computational scientometrics, science policy, and information systems. Earlier studies focused primarily on univariate time-series extrapolation, Gompertz growth models, and S-curve fitting of publication counts. More recently, computational approaches have integrated machine learning regressors and neural networks to model non-linear scientific expansion \citep{liang2021}.

\section{Bibliometric Indicators in Emerging Topic Prediction}
A variety of bibliometric indicators have been deployed to predict topic emergence:
\begin{itemize}
    \item \textbf{Volume and Growth Dynamics:} Past publication counts, acceleration rates, and momentum.
    \item \textbf{Citation Trajectories:} Early citation velocity, citation percentile ranks, and international citation diffusion.
    \item \textbf{Textual and Semantic Signals:} Keyword burstiness, term co-occurrence density, and embedding drift.
\end{itemize}
While these indicators capture intrinsic topic characteristics, they predominantly measure aggregate global signals and cannot capture localized national incubators of specialized capability.

\section{National Specialisation in Science Dynamics}
Recent empirical work has examined how countries evolve their scientific specializations over time. \citet{galanmena2026} applied network science and economic complexity methods to forecast country-level scientific diversification across 27 disciplines, demonstrating that accumulated national capabilities predict which research areas a country will enter. However, their inquiry focused on predicting the future research profile of individual countries, rather than examining whether country-level specialization predicts future global topic growth.

\section{Research Gap}
Despite the extensive literature on topic emergence and national specialization, a critical research gap persists:
\begin{enumerate}
    \item \textbf{Lack of Cross-Scale Linkage:} Existing forecasting models treat global topic growth as an endogenous function of global history, overlooking whether national specialisation in major research powers (such as Germany, China, and Japan) acts as a leading indicator of global expansion.
    \item \textbf{Inadequate Control for Database Expansion:} Many prior studies evaluate raw publication growth without adjusting for underlying bibliographic database coverage expansion, introducing systematic confounding over multi-year observation horizons.
    \item \textbf{Double-Counting in Collaborative Works:} Most comparative studies apply whole counting to internationally co-authored papers, distorting comparative advantage metrics through multi-country inflation.
\end{enumerate}
This thesis directly addresses this gap by implementing an aligned, fractionally counted, and temporally evaluated forecasting framework across 1,025 technology topics in OpenAlex.

\chapter{Methodology} \label{chap:method}

\section{OpenAlex Data}
% [Insert description of OpenAlex data source, Physical Sciences scope, and the 1,025 topics taxonomy]

\section{Fractional Counting}
% [Insert explanation of Option C: 1/k country-level equal fractional counting vs. whole counting]

\section{Dataset Construction}
% [Insert description of balanced panel structure (15,375 rows, 26 columns, 2010--2024)]

\section{Revealed Comparative Advantage (RCA)}
% [Insert formulation of Fractional All-Science RCA and denominator synchronization]

\section{Feature Engineering}
% [Insert definitions of target variable G_adj and baseline momentum M0_adj]

\section{Forecasting Models (M0--M3)}
% [Insert description of model hierarchy: M0 (Momentum), M1 (Global), M2 (Country Volumes), M3 (Specialisation)]

\section{Temporal Evaluation}
% [Insert description of test horizons (2018, 2021), origin-specific eligibility (>= 50 works), and leakage-free cutoff]

\section{Bootstrap Inference}
% [Insert evaluation metrics (MAE for RQ2, Precision@50 for RQ3) and bootstrap confidence intervals]

\chapter{Results and Discussions} \label{chap:res}

\section{RQ1: Descriptive Analysis}
% [Insert descriptive results on longitudinal trends of technology topics and specialisation in DE, CN, JP]

\section{RQ2: Forecast Accuracy}
% [Insert comparative forecasting results across M0, M1, M2, and M3 with MAE and bootstrap CIs]

\section{RQ3: Early Identification}
% [Insert results on identifying emerging high-growth topics (Top-50 precision against 20% chance level)]

\section{Discussion}
% [Insert academic and policy discussion interpreting the findings]

\chapter{Conclusions} \label{chap:conc}

\section{Conclusions}
% [Insert synthesis of main findings addressing the core thesis question]

\section{Contributions}
% [Insert summary of academic, methodological, and policy contributions]

\section{Limitations}
% [Insert discussion of research limitations: publication proxies, retrospective classifiers, country scope]

\section{Future Works}
% [Insert recommendations and avenues for future research]


%%%%%%%%%%%%%%%%%%%%%%%%%%%%%%%%%%%%%%%%%%%%%%%%%%%%%%%%%%%%%%%%%%%%
%% Appendices
%%%%%%%%%%%%%%%%%%%%%%%%%%%%%%%%%%%%%%%%%%%%%%%%%%%%%%%%%%%%%%%%%%%%
\appendix
\chapter{Independent Audit and Quality Control Verification} \label{app:audit}

To guarantee complete dataset integrity and computational reproducibility, the final primary modeling panel (\texttt{openalex\_modeling\_panel\_1025topics\_v2\_synchronized.csv}) was evaluated against an automated 29-check programmatic audit gate (\texttt{independent\_audit\_v2.py}). All 29 assertions returned a hard pass status (\texttt{EXIT\_CODE=0}).

\begin{table}[htbp]
    \centering
    \caption{Complete 29-Check Verification Gate Protocol}
    \label{tab:app_audit_checks}
    \small
    \begin{tabular}{lll}
        \toprule
        \textbf{ID} & \textbf{Validation Dimension} & \textbf{Empirical Result / Metric} \\
        \midrule
        \#01 & Master Panel v1 Preservation & SHA-256 match (\texttt{37de2b5d...}) \\
        \#02 & Original Modeling Panel v2 Preservation & SHA-256 match (\texttt{3d869014...}) \\
        \#03 & Synchronized Modeling Panel v2 Verification & SHA-256 match (\texttt{f0049cd7...}) \\
        \#04 & Total Matrix Row Count & Exactly 15,375 rows \\
        \#05 & Total Schema Column Count & Exactly 26 columns \\
        \#06 & Unique Technology Topics Universe & Exactly 1,025 unique topics \\
        \#07 & Temporal Coverage (2010--2024) & Exactly 15 annual intervals \\
        \#08 & Balanced Panel Structure & Exactly 15 observations per topic \\
        \#09 & Key Uniqueness & Zero duplicate \texttt{(topic\_id, year)} keys \\
        \#10 & Matrix Completeness & Zero null or missing values (0 / 399,750 cells) \\
        \#11 & Denominator Table Coverage & Exactly 15 annual denominator rows \\
        \#12 & Denominator Table Concordance & 100\% match with coverage census \\
        \#13 & Global All-Science Denominator & 0 mismatched rows across 15,375 rows \\
        \#14 & Germany Whole Denominator & 0 mismatched rows across 15,375 rows \\
        \#15 & Germany Fractional Denominator & 0 mismatched rows across 15,375 rows \\
        \#16 & China Whole Denominator & 0 mismatched rows across 15,375 rows \\
        \#17 & China Fractional Denominator & 0 mismatched rows across 15,375 rows \\
        \#18 & Japan Whole Denominator & 0 mismatched rows across 15,375 rows \\
        \#19 & Japan Fractional Denominator & 0 mismatched rows across 15,375 rows \\
        \#20 & Global Count Non-Negativity & Zero negative values \\
        \#21 & Germany Fractional Count Bounds & $0 \le P_{\text{Frac}} \le P_{\text{Whole}}$ (0 violations, retention: 75.35\%) \\
        \#22 & China Fractional Count Bounds & $0 \le P_{\text{Frac}} \le P_{\text{Whole}}$ (0 violations, retention: 90.03\%) \\
        \#23 & Japan Fractional Count Bounds & $0 \le P_{\text{Frac}} \le P_{\text{Whole}}$ (0 violations, retention: 87.46\%) \\
        \#24 & Germany Whole RCA Recalculation & $\Delta_{\text{max}} \le 0.00000500$ (0 exceedances) \\
        \#25 & Germany Fractional RCA Recalculation & $\Delta_{\text{max}} \le 0.00000500$ (0 exceedances) \\
        \#26 & China Whole RCA Recalculation & $\Delta_{\text{max}} \le 0.00000500$ (0 exceedances) \\
        \#27 & China Fractional RCA Recalculation & $\Delta_{\text{max}} \le 0.00000500$ (0 exceedances) \\
        \#28 & Japan Whole RCA Recalculation & $\Delta_{\text{max}} \le 0.00000500$ (0 exceedances) \\
        \#29 & Japan Fractional RCA Recalculation & $\Delta_{\text{max}} \le 0.00000500$ (0 exceedances) \\
        \bottomrule
    \end{tabular}
\end{table}

\chapter{Data Dictionary and Pilot Study} \label{app:data_dict}

\section{Master Modeling Panel Data Dictionary}
% [Insert 26-column variable definitions and schemas]

\section{Phase 1 Pilot Feasibility Study}
% [Insert documentation of Phase 1 50-topic pilot study, Option C validation, and r = 0.9968 note]


%%%%%%%%%%%%%%%%%%%%%%%%%%%%%%%%%%%%%%%%%%%%%%%%%%%%%%%%%%%%%%%%%%%%
%% Bibliography
%%%%%%%%%%%%%%%%%%%%%%%%%%%%%%%%%%%%%%%%%%%%%%%%%%%%%%%%%%%%%%%%%%%%
\clearpage
\bibliography{ref}

\end{document}
